\chapter{Thesis Plan}

\section{Working Title}
\label{sec:plan-title}

\textbf{Graph Neural Networks for Decentralized Multi-Agent Reinforcement
Learning in Power-Grid Topology Control}

\section{Central Thesis Question}
\label{sec:plan-question}

\textbf{How can graph neural networks help decentralized
reinforcement learning agents learn reliable power-grid topology control
policies that preserve survival performance while scaling to combinatorial
action spaces?}
\\
\\
Three terms fix what the thesis has to measure:

\begin{itemize}
  \item \textbf{Reliable} is episodic survival on held-out chronics, always
  reported from the checkpoint with the best test survival and re-evaluated on
  the full test split.
  \item \textbf{Decentralized} means each agent owns a fixed group of
  substations, observes a local subgraph, and acts on its own discrete topology
  actions without having any information about the other agent actions; only the critic is centralized.
  \item \textbf{Scaling} means larger grids, which a graph encoder addresses because its
  parameters are independent of grid size.
\end{itemize}

\section{Research Questions}
\label{sec:plan-research-questions}

\begin{enumerate}
  \item \textbf{Can decentralized MAPPO first learn a strong, stable
  baseline?} The flat MLP actor must achieve reliable survival on
  \texttt{bus14} before either its objective or its representation is changed.
  (Chapter~\ref{ch:baseline}.)

  {\color{blue}
  \item \textbf{Can the agents intervene sparsely without losing survival?}
  Evaluation-time rho overrides, an intervention-gated actor, fixed non-idle
  penalties, and adaptive Lagrangian intervention budgets are evaluated on two
  axes at once, survival and intervention rate, since improving either alone is
  not a result. (Chapter~\ref{ch:sparse-control}.)
  }

  \item \textbf{Can a graph actor replace the flat one without paying for it?}
  An MLP's input width is tied to one grid, whereas a shared graph encoder is
  size-independent. That advantage matters only if the GNN first matches the
  tuned MLP on \texttt{bus14}. (Chapter~\ref{ch:experiments}.)

  \item \textbf{How should the grid be represented as a graph for a local
  actor?} Node granularity, relations, substation-level augmentations, and graph
  readout are all design choices. Chapter~\ref{ch:methods} defines them and
  Chapter~\ref{ch:experiments} measures which ones move survival.

  \item \textbf{What is a local agent entitled to observe, and does enforcing
  that boundary matter?} A local graph carries contextual nodes from
  neighboring substations. Making context visible only through a live boundary
  connection, and choosing whether pooling covers the visible graph or only the
  controlled domain, determine what the policy conditions on. The distinction
  is between context and leakage. (Sections~\ref{subsec:graph-local-global}
  and~\ref{sec:graph-pooling}.)

  \item \textbf{Does any of this transfer to a larger grid?} Moving from
  \texttt{bus14} to the 36-substation WCCI grid tests whether a graph encoder is
  portable in practice and not only in parameter shape, and what has to be
  corrected: mainly feature scaling and the size of the action
  space.
  (Chapter~\ref{ch:wcci-transfer}.)
\end{enumerate}

\section{Expected Contributions}
\label{sec:plan-contributions}

\begin{itemize}
  \item \textbf{A tuned decentralized MAPPO baseline on \texttt{bus14}.} The
  original framework trained unstably; the training and initialization changes
  that fix it are isolated one at a time.

  \item \textbf{An ablation of sparse intervention mechanisms} reported jointly
  on survival and intervention frequency, connecting the learned behavior to
  what a control room would accept: act rarely, act locally, and act when the
  grid is unsafe.

  \item \textbf{A shared GNN actor and a screened local graph design.} The GNN
  first matches the tuned MLP, then three representations are compared using a
  fixed candidate graph with a dynamic edge mask and an explicit information
  boundary. Encoder depth and width, message-passing operator, readout, and
  structural augmentations are screened with full-test evaluation.

  \item \textbf{A size-invariant mechanism for combinatorial action spaces.} A
  candidate-action head that scores each action from the encoded nodes it
  modifies, replacing the fixed-width layer over an action index.

  \item \textbf{A measurement of what actually breaks under transfer.} Encoder
  parameters cross grids unchanged; the input distribution they were calibrated
  on does not. Two findings generalize beyond power grids: per-feature
  standardization is harmful when a channel is a mixture dominated by structural
  zeros, and for gauge-dependent quantities such as voltage angle, changing the
  representation beats any choice of normalization.
\end{itemize}

\section{Open Decisions}

\begin{itemize}
  \item Confirm the final thesis title and whether to emphasize GNNs,
  sparse-control, or decentralized MARL most strongly.
  \item Decide which experiment families are final and which are only
  exploratory.
  \item Decide whether the main claim is performance improvement, improved
  sparsity at equal survival, or reproducible ablation insight.
  \item Choose the final bibliography style required by the university.
\end{itemize}
